\section{Code Generation Experiment (MBPP)}
\label{app:mbpp}

To validate that CJE generalizes beyond preference judgments, we evaluated on MBPP code generation \citep{Austin2021MBPP}, where ground truth is objective: unit tests either pass or fail.

\paragraph{Setup.}
We use 257 MBPP problems (974 samples across 5 epochs) with five policies: \texttt{base} (GPT-4o standard, 87.4\% pass rate, used for calibration), \texttt{enhanced} (GPT-4o with expert prompting, 87.2\%), \texttt{terse} (GPT-4o with minimal prompting, 87.4\%), \texttt{mini} (GPT-4o-mini, 80.9\%), and \texttt{unhelpful} (deliberately buggy, 56.0\%).
The judge is GPT-4o-mini predicting pass probability (0--100, normalized to 0--1).
Oracle labels are actual unit test results.
Unlike Arena (where response length confounds judge scores), MBPP has no strong covariate effects, so we use monotone-only calibration.

\paragraph{Transport validation.}
We apply the mean transport test (\cref{prop:mean-transport}) to each target policy. Two policies fail:
\begin{itemize}[leftmargin=*,nosep]
\item \texttt{mini}: Mean residual $-0.065$ ($p{=}0.006$). The judge model (GPT-4o-mini) is identical to this policy, creating a judge--policy confound that violates transport.
\item \texttt{unhelpful}: Mean residual $-0.280$ ($p{<}0.001$). Low-support saturation: 31\% of samples score below base's 10th percentile, where calibrator training density is sparse.
\end{itemize}
Both policies are excluded from RMSE computation; including them would bias results.

\paragraph{Results.}
\Cref{tab:mbpp-results} shows CJE performance across oracle fractions for the two transport-valid policies (\texttt{enhanced}, \texttt{terse}) at full sample size ($n{=}974$).
Calibration provides consistent RMSE reduction across all oracle fractions: even at 5\% oracle (${\sim}50$ labels), CJE achieves 18\% improvement over naive estimation (raw judge scores without calibration).
At 10\%+ oracle, improvements reach 32--84\%.

\begin{table}[h]
\centering
\small
\caption{\textbf{CJE on MBPP code generation: oracle fraction tradeoff.}
RMSE computed over two transport-valid policies (\texttt{enhanced}, \texttt{terse}) at $n{=}974$.
Unlike Arena where 5\% oracle suffices for ranking, MBPP shows diminishing returns: 10\% oracle captures most of the benefit.}
\label{tab:mbpp-results}
\begin{tabular}{lcccc}
\toprule
Oracle \% & \# Labels & CJE & Naive (uncalib.) & Improvement \\
\midrule
5\%   & 49   & 0.039 & 0.048 & +18\% \\
10\%  & 97   & 0.033 & 0.048 & +32\% \\
25\%  & 244  & 0.019 & 0.048 & +60\% \\
50\%  & 487  & 0.014 & 0.048 & +70\% \\
100\% & 974  & 0.008 & 0.048 & +84\% \\
\bottomrule
\end{tabular}
\end{table}

\paragraph{Comparison to Arena.}
Both experiments validate the oracle fraction tradeoff, but with different thresholds:
\begin{itemize}[leftmargin=*,nosep]
\item \emph{Arena}: 5\% oracle achieves 94\% ranking accuracy; calibration benefit plateaus early.
\item \emph{MBPP}: 5\% oracle provides 18\% RMSE reduction; 10\% provides 32\%; benefit continues scaling to 84\% at full oracle.
\end{itemize}
The difference likely reflects task structure: Arena evaluates preferences (ordinal, high variance) while MBPP evaluates correctness (binary, low variance). Binary oracles require more labels to stabilize calibration.

\paragraph{Judge--policy confound.}
The \texttt{mini} transport failure highlights a methodological risk: when a model judges its own outputs, self-preference or style-alignment biases can corrupt calibration.
Practitioners should use external judges or validate transport explicitly.

\paragraph{Takeaway.}
CJE successfully transfers to code generation with objective ground truth.
The oracle fraction tradeoff generalizes: more labels enable larger gains.
Transport validation is essential---without it, confounded policies inflate RMSE by 2--3$\times$.
